\slajdid{09_2-s039}
\begin{frame}[label=09_2-s039]
\fontsize{10}{12}\selectfont\setlength{\abovedisplayskip}{5pt}\setlength{\belowdisplayskip}{5pt}\setlength{\belowdisplayshortskip}{3pt}\setlength{\abovedisplayshortskip}{2pt}
Tranzistori sa dugačkim kanalom $V_S\ll E_CL$\par\medskip
U slučaju kada su tranzistori sa dugačkim kanalom, odnosno zanemaren je efekat zasićenja brzine nosilaca izraz se svodi na
\[k_p(V_I-V_{DD}-V_{Tp})^2=k_n(V_I-V_{Tn})^2\]
odnosno
\[-\sqrt{\frac{k_p}{k_n}}(V_S-V_{DD}-V_{Tp})=(V_S-V_{Tn})\]
(pazite na oslobađanje kvadrata u ovakvim situacijama – leva strana mora da bude pozitivna)\par\medskip
Prilikom ovih izvođenja česta situacija u literaturi jeste da se uvede smena
\[\sqrt{\frac{k_p}{k_n}}=r\]
r – odnos (ratio), pa je
\[V_S=\frac{V_{Tn}+r(V_{DD}+V_{Tp})}{1+r}\]
\end{frame}
